\section{Background and motivation}\label{sec:background}

\subsection{Near-UV light delivery in trapped-ion experiments}
In \Ybion{} experiments, light at \SI{369.5}{nm} drives the \ce{^2S_{1/2}}--\ce{^2P_{1/2}} transition used for
Doppler cooling, optical pumping and state detection. Light at \SI{399}{nm} is used for isotope-selective
photoionisation of neutral Yb. Delivering this light through single-mode fibres has clear advantages.
The fibre acts as a spatial filter that produces a clean Gaussian mode. It decouples the laser table
from the vacuum system and keeps beam pointing at the ion stable. With polarisation-maintaining (PM)
fibre, it also delivers a well-defined polarisation. Scalable architectures with many ion-trap modules
depend even more on fibre delivery, so the long-term reliability of near-UV fibres is a practical
concern for the whole field.

\subsection{The problem in the group}
The group uses Thorlabs PM patch cables designed for the near UV (referred to in the lab as
``375PM''; \cref{sec:sota-fibres}) at powers of a few milliwatts. The behaviour observed is:
\begin{itemize}
  \item directly after the end faces have been polished, the coupling efficiency is excellent;
  \item it then declines over time, faster when more power is used;
  \item both the input and the output end are affected;
  \item polishing the end faces again restores the throughput.
\end{itemize}
The practical consequences are frequent re-polishing (person-hours and risk of damaging the connector),
slow drifts of the power delivered to the ions, the need for spare laser power, and occasionally
sudden failures during measurements.

\paragraph{What these observations already tell us} Re-polishing removes only micrometres to tens of
micrometres of glass. Colour centres formed along the length of the fibre would survive it. The
restoration therefore places the damage at, or within micrometres of, the end faces. The power dependence
shows that the process is driven by light. The degradation of the \emph{output} face, which sees almost the
same guided intensity as the input face but no focusing optics, argues against simple misalignment of the
input coupling and for a process driven by local intensity at a glass--gas interface. At \SI{3}{mW} in a
\SI{2.3}{\micro m} mode-field diameter, the peak intensity is \SI{\approx1.4e5}{W/cm^2}, and each
square centimetre of the core region receives about \num{3e23} photons per second (\cref{app:estimates}).
At such photon fluxes even very weak absorption by surface sites or adsorbed molecules leads to
photochemistry within seconds.

\subsection{Terminology and scope}
In the lab, ``solarisation'' is used for any UV-induced loss of fibre throughput. Strictly, solarisation
means the formation of colour centres in the glass (\cref{sec:sota-bulk}). In this document,
\emph{end-face degradation} denotes the observed phenomenon and \emph{solarisation} denotes colour-centre
formation. Finding out which of the two is happening, or whether both are, is part of the project (RQ1).
The scope is the atmosphere at the end faces of the fibre type used in the group, at \SI{369.5}{nm} and
a nearby testbed wavelength (\SI{375}{nm}). Developing new fibres is out of scope, but alternative fibre
technologies are reviewed (\cref{sec:sota-fibres}) because they are the fall-back if the atmosphere turns
out not to be the cause.
